% =============================================================================
% PR\'EFACE — Apprentissage par Renforcement Profond
% =============================================================================
\chapter*{Pr\'eface}
\addcontentsline{toc}{chapter}{Pr\'eface}
\markboth{Pr\'eface}{Pr\'eface}

\section*{Objectifs de cet ouvrage}

L'apprentissage par renforcement profond (\emph{Deep Reinforcement Learning}, DRL)
combine l'apprentissage par renforcement classique et l'apprentissage profond.
Au cours de la derni\`ere d\'ecennie, cette discipline a permis des avanc\'ees
spectaculaires~: vaincre des champions humains aux jeux de Go et d'Atari,
contr\^oler des robots manipulateurs avec une dext\'erit\'e remarquable, optimiser
des strat\'egies financi\`eres et concevoir de nouvelles mol\'ecules th\'erapeutiques.

Cet ouvrage s'adresse aux \'etudiants de Master et de Doctorat en informatique,
math\'ematiques appliqu\'ees ou sciences de l'ing\'enieur, ainsi qu'aux chercheurs et
praticiens souhaitant acqu\'erir une compr\'ehension rigoureuse des fondements
th\'eoriques et des m\'ethodes algorithmiques du DRL.

\section*{Organisation du cours}

Le cours est structur\'e en onze chapitres, organis\'es selon une progression p\'edagogique
allant des fondements math\'ematiques aux applications de pointe~:

\begin{enumerate}
    \item \textbf{Processus de d\'ecision markoviens (MDP)}~: le cadre math\'ematique
    fondamental qui formalise l'interaction agent-environnement.

    \item \textbf{Programmation dynamique}~: les algorithmes d'it\'eration sur la valeur
    et d'it\'eration sur la politique, fondements de toute m\'ethode de RL.

    \item \textbf{M\'ethodes de Monte Carlo et diff\'erences temporelles}~: les approches
    sans mod\`ele fond\'ees sur l'\'echantillonnage.

    \item \textbf{Q-Learning et SARSA}~: les algorithmes tabulaires fondamentaux
    d'apprentissage hors-politique et en-politique.

    \item \textbf{R\'eseaux Q profonds (DQN)}~: la r\'evolution de l'approximation de
    fonction par r\'eseaux de neurones appliqu\'ee au RL.

    \item \textbf{Gradient de politique --- REINFORCE}~: l'optimisation directe de la
    politique par mont\'ee de gradient.

    \item \textbf{M\'ethodes acteur-critique (A3C, PPO)}~: la combinaison de
    l'estimation de la valeur et de l'optimisation de la politique.

    \item \textbf{\'Etat de l'art (SAC, TD3, DDPG)}~: les algorithmes modernes
    pour les espaces d'action continus.

    \item \textbf{Apprentissage par renforcement multi-agents}~: la g\'en\'eralisation
    aux syst\`emes \`a agents multiples.

    \item \textbf{RL sous contraintes et s\^uret\'e}~: la prise en compte de contraintes
    de s\'ecurit\'e dans l'apprentissage.

    \item \textbf{Applications}~: robotique, jeux, finance et au-del\`a.
\end{enumerate}

\section*{Pr\'erequis}

Le lecteur devra poss\'eder des connaissances solides en~:

\begin{itemize}
    \item \textbf{Probabilit\'es et statistiques}~: variables al\'eatoires, esp\'erance
    conditionnelle, cha\^ines de Markov, convergence stochastique.

    \item \textbf{Optimisation}~: descente de gradient, convexit\'e, multiplicateurs de
    Lagrange, m\'ethodes num\'eriques.

    \item \textbf{Alg\`ebre lin\'eaire}~: espaces vectoriels, d\'ecompositions matricielles,
    valeurs propres.

    \item \textbf{Apprentissage automatique}~: r\'eseaux de neurones, r\'etropropagation,
    r\'egularisation, architectures profondes (CNN, RNN).

    \item \textbf{Programmation Python}~: ma\^itrise de NumPy, PyTorch et des
    environnements Gymnasium (anciennement OpenAI Gym).
\end{itemize}

\section*{Notations}

Nous adoptons les conventions suivantes tout au long de l'ouvrage~:

\begin{center}
\renewcommand{\arraystretch}{1.3}
\begin{tabular}{cl}
    \toprule
    \textbf{Symbole} & \textbf{Signification} \\
    \midrule
    $\mathcal{S}$ & Ensemble des \'etats \\
    $\mathcal{A}$ & Ensemble des actions \\
    $\mathcal{R}$ & Ensemble des r\'ecompenses \\
    $\gamma \in [0,1)$ & Facteur d'actualisation \\
    $\pi(a \mid s)$ & Politique stochastique \\
    $V^\pi(s)$ & Fonction de valeur d'\'etat sous $\pi$ \\
    $Q^\pi(s,a)$ & Fonction de valeur d'action sous $\pi$ \\
    $A^\pi(s,a)$ & Fonction avantage sous $\pi$ \\
    $P(s' \mid s, a)$ & Probabilit\'e de transition \\
    $R(s, a, s')$ & Fonction de r\'ecompense \\
    $\E[\cdot]$ & Esp\'erance math\'ematique \\
    $\PP(\cdot)$ & Probabilit\'e \\
    $\nabla_\theta$ & Gradient par rapport \`a $\theta$ \\
    $\theta$ & Param\`etres de la politique ou du r\'eseau \\
    $\alpha$ & Taux d'apprentissage \\
    $\epsilon$ & Param\`etre d'exploration ($\epsilon$-greedy) \\
    $\tau$ & Coefficient de mise \`a jour douce \\
    $\mathcal{D}$ & M\'emoire de rejeu (\emph{replay buffer}) \\
    \bottomrule
\end{tabular}
\end{center}

\section*{Conventions typographiques}

\begin{itemize}
    \item Les \textbf{d\'efinitions} sont pr\'esent\'ees dans des encadr\'es bleus.
    \item Les \textbf{th\'eor\`emes} et \textbf{propositions} apparaissent dans des
    encadr\'es formels avec preuves.
    \item Les \textbf{algorithmes} sont pr\'esent\'es en pseudo-code dans des encadr\'es
    d\'edi\'es.
    \item Les \textbf{impl\'ementations} PyTorch/Gymnasium sont donn\'ees dans des blocs
    de code syntaxiquement color\'es.
    \item Les \textbf{formules cl\'es} sont mises en \'evidence dans des encadr\'es
    sp\'eciaux.
    \item Les points n\'ecessitant une \textbf{attention particuli\`ere} sont signal\'es
    dans des encadr\'es d'avertissement.
\end{itemize}

\section*{Remerciements}

Cet ouvrage est le fruit de plusieurs ann\'ees d'enseignement et de recherche en
apprentissage par renforcement. Nous remercions les \'etudiants qui, par leurs questions
et leur curiosit\'e, ont contribu\'e \`a am\'eliorer la pr\'esentation de ces sujets.
Nous exprimons \'egalement notre gratitude envers la communaut\'e open source, en
particulier les d\'eveloppeurs de PyTorch, Gymnasium, Stable-Baselines3 et CleanRL,
dont les outils rendent l'apprentissage par renforcement accessible \`a tous.

\section*{Comment utiliser ce livre}

\begin{intuition}
Chaque chapitre suit une structure coh\'erente~:
\begin{enumerate}
    \item \textbf{Motivation et intuition} --- Pourquoi cette m\'ethode~? Quel probl\`eme
    r\'esout-elle~?
    \item \textbf{Fondements th\'eoriques} --- D\'efinitions, th\'eor\`emes, preuves de
    convergence.
    \item \textbf{Algorithmes} --- Pseudo-code d\'etaill\'e avec analyse de complexit\'e.
    \item \textbf{Impl\'ementation} --- Code PyTorch fonctionnel et reproductible.
    \item \textbf{Exercices} --- Probl\`emes th\'eoriques et pratiques de difficult\'e
    croissante.
\end{enumerate}
\end{intuition}

\section*{Bref historique}

L'apprentissage par renforcement plonge ses racines dans les travaux de Bellman
(programmation dynamique, ann\'ees 1950) et de la psychologie comportementale
(conditionnement op\'erant de Skinner). Les jalons majeurs incluent~:

\begin{itemize}
    \item \textbf{1989}~: Watkins introduit le Q-Learning, premier algorithme
    hors-politique convergent pour les MDP.
    \item \textbf{1992}~: TD-Gammon de Tesauro apprend le backgammon par auto-jeu
    avec diff\'erences temporelles.
    \item \textbf{2013}~: Mnih et al.\ (DeepMind) combinent DQN avec des r\'eseaux
    convolutifs pour jouer aux jeux Atari \`a partir de pixels.
    \item \textbf{2015}~: Publication de DQN dans \emph{Nature}~; Silver et al.\ d\'eveloppent
    AlphaGo.
    \item \textbf{2016}~: AlphaGo bat Lee Sedol, champion mondial de Go.
    \item \textbf{2017}~: Schulman et al.\ introduisent PPO~; Haarnoja et al.\ proposent SAC.
    \item \textbf{2019}~: OpenAI Five bat les champions du monde de Dota~2.
    \item \textbf{2020--2024}~: RL appliqu\'e \`a la d\'ecouverte de m\'edicaments, \`a la
    fusion nucl\'eaire (DeepMind/EPFL), \`a l'alignement de LLM (RLHF).
\end{itemize}

\section*{Ressources compl\'ementaires}

Pour approfondir les sujets abord\'es dans cet ouvrage, nous recommandons~:

\begin{itemize}
    \item \textbf{Sutton \& Barto} --- \emph{Reinforcement Learning: An Introduction}
    (2\textsuperscript{e} \'edition, 2018). L'ouvrage de r\'ef\'erence en RL classique.

    \item \textbf{Bertsekas} --- \emph{Dynamic Programming and Optimal Control}
    (4\textsuperscript{e} \'edition). Traitement rigoureux de la programmation dynamique.

    \item \textbf{Szepesv\'ari} --- \emph{Algorithms for Reinforcement Learning}.
    Synth\`ese concise des algorithmes fondamentaux.

    \item \textbf{Agarwal et al.} --- \emph{Reinforcement Learning: Theory and Algorithms}
    (2022). Traitement th\'eorique moderne.

    \item \textbf{Cours en ligne}~: David Silver (UCL), Sergey Levine (UC Berkeley),
    Emma Brunskill (Stanford).
\end{itemize}

\vfill

\begin{flushright}
\textit{[Author Name]} \\
\textit{Mars 2026}
\end{flushright}
